% questions.tex —— 1958 年·安徽卷（存疑） · 题目内容片段（无答案）
% 由 tools/generate_projects.py 依来源总账自 1958/1958_安徽_数学真题_试题.md 生成，勿手改（改源 MD 后重跑）


\begin{gktip}
\textbf{来源说明}：本卷题面出自登记版固定字节文件，上游与文件名均标注「来源网络，无法确认真实性」——本卷来源存疑，转录如实照录，真实性未能核实。原卷各题小问以拉丁字母 (a)(b)(c)… 编号，本转录按体例统一为（1）（2）（3）…，内容不变。
\end{gktip}

\begin{question}
\gkti{1}（每小题7分，计35分）

（1） 已知方程 $x^2-2x+q=0$ 两根的差的平方等于16，求 $q$ 的值。

（2） 计算 $\log_{\frac32}\dfrac94+(0.25)^{-\frac32}+\log_2\sin135°$。

（3） 证明 $\sin(\alpha+\beta)<\sin\alpha+\sin\beta$（$\alpha,\beta$ 都是锐角）。

（4） 说明函数 $y=x^2-4$ 的图象和函数 $y=x^2$ 的图象对于坐标轴的位置有什么区别。

（5） 夹在两个平行平面间的两条线段，它们的长的和是28寸，它们在同一个平面上的射影是9寸和5寸，求这两条线段的长。

\end{question}
\begin{question}
\gkti{2}（每小题10分，计20分）

\Needspace{5\baselineskip}
（1） 解方程组

\nopagebreak[4]
\[
\begin{cases}
\log_3(2x+8)-\log_3 y=1,\\
\sqrt{2x+14}-\sqrt{2y+7}=0.
\end{cases}
\]

（2） 解方程 $\left(\sqrt{2+\sqrt3}\right)^x+\left(\sqrt{2-\sqrt3}\right)^x=4$。

\end{question}
\begin{question}
\gkti{3}（14分）设 $\sin\alpha$ 为 $\sin\theta$ 与 $\cos\theta$ 的等差中项，$\sin\beta$ 为 $\sin\theta$ 与 $\cos\theta$ 的等比中项，求证：

\nopagebreak[4]
\[
2\cos2\alpha=\cos2\beta
\]

\end{question}
\begin{question}
\gkti{4}（15分）已知三棱锥的底面积为 $S$，高为 $H$，过这个三棱锥的三个侧棱的中点作一截面，求过这个三棱锥的底面上任一点为顶点，以所作截面为底的新的三棱锥的体积。

\end{question}
\begin{question}
\gkti{5}（16分）在内接于圆的四边形 $ABCD$ 中，对应边 $AB$、$DC$ 的延长线交于 $P$，$AD$、$BC$ 的延长线交于 $Q$，求证：

（1） $PB\cdot CA=PC\cdot BD$；

（2） 角 $P$ 及角 $Q$ 的平分线互相垂直。

\end{question}